\documentclass[british]{article}

\usepackage[margin=1.5in]{geometry}

\title{Automorphisms of holomorphic foliations on rational surfaces}
\author{Percy Fernández Sánchez}
\date{2002}

% Standard
\usepackage{amssymb,amsmath,amscd}
\usepackage{enumerate}
\usepackage{tikz-cd}
\usepackage{mathtools}
\usepackage{booktabs}


%Bibliography
\usepackage[backend=biber,maxnames=99,maxalphanames=4]{biblatex}
\addbibresource{PM-16-2002-47.bib}


% Colours
\usepackage{xcolor}


% Fonts
\usepackage{mathrsfs}
\usepackage{fouriernc}


% Things for Quarto/pandoc
\definecolor{quarto-callout-color-frame}{HTML}{000000}
\usepackage{calc,array}
\usepackage{longtable}
\providecommand{\tightlist}{%
\setlength{\itemsep}{0pt}\setlength{\parskip}{0pt}}
\newcounter{none}
\usepackage{graphicx}
\makeatletter
\newsavebox\pandoc@box
\newcommand*\pandocbounded[1]{% scales image to fit in text height/width
  \sbox\pandoc@box{#1}%
  \Gscale@div\@tempa{\textheight}{\dimexpr\ht\pandoc@box+\dp\pandoc@box\relax}%
  \Gscale@div\@tempb{\linewidth}{\wd\pandoc@box}%
  \ifdim\@tempb\p@<\@tempa\p@\let\@tempa\@tempb\fi% select the smaller of both
  \ifdim\@tempa\p@<\p@\scalebox{\@tempa}{\usebox\pandoc@box}%
  \else\usebox{\pandoc@box}%
  \fi%
}
% Set default figure placement to htbp
\def\fps@figure{htbp}
\makeatother
\usepackage{kvsetkeys}


% Header and footer
\usepackage{fancyhdr}
\usepackage{lastpage}
\usepackage{datetime2}
\pagestyle{fancy}
\fancypagestyle{plain}{}
\fancyhf{}
\renewcommand{\headrulewidth}{0pt}
\cfoot{\small\thepage\ of \pageref*{LastPage}}
\lfoot{\footnotesize Build date: \today}


% Theorem environments
\usepackage{amsthm}
\newenvironment{itenv}[1]
  {\phantomsection\par\smallskip\noindent\textbf{#1.}\itshape}
  {\par\smallskip}
\newenvironment{rmenv}[1]
  {\phantomsection\par\smallskip\noindent\textbf{#1.}\rmfamily}
  {\par\smallskip}


% Shortcuts
\newcommand{\oldpage}[1]{\marginpar{\footnotesize$\Big\vert$ \textit{p.~#1}}}


%% Document %%

% hyperref last, as is "good practice"

\usepackage{hyperref}
\hypersetup{colorlinks,linkcolor={purple!50!black},citecolor={purple!50!black},urlcolor={purple!80!black}}

\usepackage{embedall}
\begin{document}

\maketitle
\thispagestyle{fancy}

\setcounter{footnote}{0}


%% Content %%

\emph{This document is a translation into English of the following:}

\begin{quote}
Sánchez, PF. ``Automorfismo de foliaciones holomorfas sobre superficies
racionales''. \emph{Pro Mathematica} \textbf{16} (2002), 47--59.
\href{https://revistas.pucp.edu.pe/index.php/promathematica/article/view/8182/8478}{revistas.pucp.edu.pe/index.php/promathematica/article/view/8182/8478}.
\end{quote}

\emph{The translator (\href{https://thosgood.net}{Tim Hosgood}) takes
full responsibility for any errors introduced in this document, and
claims no rights to any of the mathematical content.}

\tableofcontents

\providecommand{\scr}[1]{{\mathscr{#1}}}
\renewcommand{\cal}[1]{{\mathcal{#1}}}
\renewcommand{\frak}[1]{{\mathfrak{#1}}}
\renewcommand{\geq}{\geqslant}
\renewcommand{\leq}{\leqslant}

\providecommand{\ZZ}{\mathbb{Z}}
\providecommand{\QQ}{\mathbb{Q}}
\providecommand{\CC}{\mathbb{C}}
\providecommand{\PP}{\mathbb{P}}
\providecommand{\id}{\mathrm{id}}
\providecommand{\dd}{\mathrm{d}}

\providecommand{\Sing}{\operatorname{Sing}}
\providecommand{\tang}{\operatorname{tang}}
\providecommand{\Aut}{\operatorname{Aut}}
\providecommand{\CS}{\operatorname{CS}}
\providecommand{\Kod}{\operatorname{Kod}}
\providecommand{\PGL}{\operatorname{PGL}}
\providecommand{\GL}{\operatorname{GL}}
\providecommand{\PSL}{\operatorname{PSL}}
\providecommand{\gl}{\operatorname{gl}}

\protect\phantomsection\label{abstract}
\begin{rmenv}{Abstract}
In this work, we classify the holomorphic foliations with infinite
automorphism group on a rational surface. As a consequence of this
result, we prove that the automorphism group of a foliation of general
type with singularities on a rational surface is finite.

\end{rmenv}

\section{Introduction}\label{introduction}

\oldpage{48}

Schwarz proved that the automorphism group of a Riemann surface of genus
greater than 2 is finite. Andreotti, in \autocite{1}, generalised this
result, proving that the group of bimeromorphisms of an algebraic
variety of general type is finite, these being analogous to Riemann
surfaces of genus greater than 2. In the case of algebraic surfaces,
there is even a bound for this number, see \autocite{14}.

In this work we first classify the holomorphic foliations with
singularities on rational surfaces.

\protect\phantomsection\label{theorem-A}
\begin{itenv}{Theorem A}
Let \({\mathcal{F}}\) be a foliation on a rational surface \(M\). If
\(\#\operatorname{Aut}({\mathcal{F}})=+\infty\) then \({\mathcal{F}}\)
is bimeromorphic to a Riccati foliation or to a rational fibration.

\end{itenv}

We then prove a result analogous to that of Andreotti for holomorphic
foliations on rational surfaces (surfaces bimeromorphic to the
projective plane).

\protect\phantomsection\label{theorem-B}
\begin{itenv}{Theorem B}
The automorphism group of a holomorphic foliation of general type with
singularities on a rational surface is finite.

\end{itenv}

\section{Preliminaries}\label{preliminaries}

For the basic notions of holomorphic foliations, we recommend the books
\autocite{4}, \autocite{13}, and \autocite{2}. A holomorphic foliation
\({\mathcal{F}}\) with isolated singularities on an algebraic surface
\(M\) can be defined by a family of holomorphic vector fields
\(\{X_i\}\) defined on an open cover \(\{U_i\}\) of \(M\) satisfying the
cocycle condition \(X_i=g_{ij}X_j\) whenever
\(U_i\cap U_j\neq\varnothing\), where the \(\{g_{ij}\}\) are
nowhere-zero holomorphic functions defined on the \(U_i\cap U_j\). In
this case, the \(\{g_{ij}\}\) define a line bundle \(T_{\mathcal{G}}^*\)
on \(M\) called the \emph{cotangent}, or \emph{canonical},
\emph{bundle}. The set of singularities
\(\operatorname{Sing}({\mathcal{G}})\) of \({\mathcal{G}}\) is defined
as \(\operatorname{Sing}({\mathcal{G}})/_{U_i}=\{X_i=0\}\), and it is
always possible to suppose that
\(\dim\operatorname{Sing}({\mathcal{G}})=0\).

Let \({\mathcal{F}}\) be a holomorphic foliation on a complex surface
\(M\). Consider a compact curve \(C\) on \(M\). \oldpage{49} Let
\(p\in C\) and let \(\{f=0\}\) be a reduced local equation of \(C\)
around \(p\). Suppose that \({\mathcal{F}}\) is represented in a
coordinate neighbourhood \((U,(x,y))\) of \(p=(0,0)\) by the holomorphic
\(1\)-form \[
  \omega
  = a(x,y)\mathrm{d}{x} + b(x,y) \mathrm{d}{y}.
\] When \(C\) is not \({\mathcal{F}}\)-invariant, we define the
\emph{tangency} between \({\mathcal{F}}\) and \(C\) at \(p\) by \[
  \operatorname{tang}_p({\mathcal{F}},C)=\dim_\mathbb{C}{\mathcal{O}}_p/I
\] where \(I\) is the ideal generated by \(f\) and
\(-b\frac{\partial f}{\partial x}+a\frac{\partial f}{\partial y}\).
Define
\(\operatorname{tang}({\mathcal{F}},C)=\sum_{p\in C}\operatorname{tang}_p({\mathcal{F}},C)\).
It is proven in \autocite{2} that \[
  T_{\mathcal{F}}^*C
  = \operatorname{tang}({\mathcal{F}},C) - C^2.
\]

\protect\phantomsection\label{example-1}
\begin{rmenv}{Example 1}
If \(M=\mathbb{C}P(2)\) and \(C\) is a non-\({\mathcal{F}}\)-invariant
line, then we have \[
  T_{\mathcal{F}}^*
  = {\mathcal{O}}_{\mathbb{C}P(2)}(\operatorname{tang}({\mathcal{F}},C)-1).
\]

\end{rmenv}

Now suppose that \(C\) is \({\mathcal{F}}\)-invariant, so that
\(\omega\wedge\mathrm{d}{f}=f\Theta\), where \(\Theta\) is a holomorphic
\(2\)-form. Then there exist relatively prime holomorphic functions
\(g\) and \(h\) defined on \(U\), along with a holomorphic \(1\)-form
\(\eta\), such that \[
  g\omega
  = h\mathrm{d}{f} + f\eta.
\]

We define the \emph{Camacho--Sad index} at \(p\) by \[
  \operatorname{CS}_p({\mathcal{F}},C)
  = \frac{-1}{2\pi i}\int_\gamma\frac{\eta}{h}
\] where \(\gamma\) is a loop around \(p\) on \(\{f=0\}\). Define
\(\operatorname{CS}({\mathcal{F}},C)=\sum_{p\in\operatorname{Sing}{\mathcal{F}}\cap C}\operatorname{CS}_p({\mathcal{F}},C)\).
The \emph{Camacho--Sad index theorem} \autocite{5} says that \[
  \operatorname{CS}({\mathcal{F}},C)
  = C^2.
\] Recall that a \emph{reduced foliation} is a foliation
\({\mathcal{F}}\) such that every singularity \(p\) is reduced in the
sense of Seidenberg, i.e.~for every vector field \(X\) that generates
\({\mathcal{F}}\), and every singular point \(p\) of \(X\), the
eigenvalues of the linear part of \(X\) are not both zero, and their
quotient, when it is defined, is not a positive rational number. If one
of the values is zero and the other non-zero, then the singularity is
said to be a \emph{saddle node}.

\oldpage{50}

\protect\phantomsection\label{definition-1}
\begin{rmenv}{Definition 1}
Let \({\mathcal{F}}\) be a foliation on the complete surface \(S\), and
\({\mathcal{G}}\) an arbitrary reduced foliation that is
bimeromorphically equivalent to \({\mathcal{F}}\). The \emph{Kodaira
dimension} of \({\mathcal{F}}\) is given by \[
  \operatorname{Kod}({\mathcal{F}})
  = \limsup_{n\to\infty} \frac{\log h^0(S,K_{\mathcal{G}}^{\otimes n})}{\log n}.
\]

\end{rmenv}

The concept of Kodaira dimension for holomorphic foliations was
introduced independently by L.G. Mendes and M. McQuillan, see
\autocite{2}, \autocite{12}, and \autocite{11}.

It has been proven that the Kodaira dimension is well defined and is a
bimeromorphic invariant of \({\mathcal{F}}\), see \autocite{12}.

When the foliation has Kodaira dimension \(2\), we say that the
foliation is of \emph{general type}. This terminology is justified since
there exists a partial classification of foliations of Kodaira dimension
less than \(2\). For more details, see \autocite{11} and \autocite{2}.

A \emph{fibration} at \(M\) on a compact Riemann surface \(B\) is a
surjective holomorphic map \(f\colon M\to B\). A \emph{generic} fibre of
\(f\) is the preimage of a regular value of \(f\). If all the generic
fibres are connected then the fibration \(f\) is said to be
\emph{connected}. If all the fibres are rational curves then the
fibration \(f\) is said to be a \emph{rational} fibration.

A foliation \({\mathcal{F}}\) at \(M\) is said to be a \emph{Riccati}
foliation if there exists a rational fibration whose fibres are
transversal to \({\mathcal{F}}\).

In \autocite{12}, it was proven that Riccati foliations and rational
fibrations have dimension at most \(1\).

The \emph{bimeromorphism} (resp. \emph{automorphism}) \emph{group} of
the foliation \({\mathcal{F}}\) is the maximal subgroup of the
bimeromorphisms (resp. automorphisms) of \(M\) that leave the foliation
invariant.

In the case of \(\mathbb{P}_\mathbb{C}^2\), a foliation
\({\mathcal{F}}\) can be defined in affine coordinates \(z=1\) by a
\(1\)-form \(\omega=A\mathrm{d}{x}+B\mathrm{d}{y}\), where \(A\) and
\(B\) are complex polynomials in the variables \(x,y\). In this case, we
can show that:

\oldpage{51}

\protect\phantomsection\label{proposition-1}
\begin{itenv}{Proposition 1}
The automorphism group of a holomorphic foliation on
\(\mathbb{P}_\mathbb{C}^2\) is a Lie group. Furthermore, it is
quasi-projective, and thus has a finite number of connected components.

\end{itenv}

\begin{proof}
By definition, \(\operatorname{Aut}({\mathcal{F}})\) is the set of
elements \(g\) of \(\operatorname{PGL}(3,\mathbb{C})\) that leave
\({\mathcal{F}}\) invariant. So if \({\mathcal{F}}\) is given by a
homogeneous \(1\)-form \(\omega=\sum P_i\mathrm{d}{x}_i\) then, by
definition, \(\omega\) and \(g^*\omega\) define the same foliation, and
are thus linearly dependent. Thus \[
  \operatorname{Aut}(cal{F})
  = \mathbb{P}(\{ g\in\operatorname{GL}(3,\mathbb{C}) \mid g^*\omega\wedge\omega = 0 \})
\] where \(\mathbb{P}(A)\) denotes the projectivisation of \(A\). It is
clear that \(\operatorname{Aut}({\mathcal{F}})\) is a Lie group. Since
\(\operatorname{GL}(3,\mathbb{C})\) is quasi-projective, the proof is
complete.
\end{proof}

\protect\phantomsection\label{example-2}
\begin{rmenv}{Example 2}
Suppose that a foliation \({\mathcal{F}}\) is defined on
\(\mathbb{P}_\mathbb{C}^2\) by the \(1\)-form
\(\omega=P(x,y)\mathrm{d}{x}+(Q(x,y)-yP(x,y))\mathrm{d}{y}\), where
\(P,Q\in\mathbb{C}[x+\frac{1}{2\alpha}y-\frac{1}{2}y^2]\) for some
\(\alpha\in\mathbb{C}^*\). Such foliations are invariant under
\(T(x,y)=(x+\frac{1}{\alpha}y,y+\frac{1}{\alpha})\). So
\(\operatorname{Aut}({\mathcal{F}})\) contains infinitely many elements
of \(\{T^n\mid n\in\mathbb{Z}\}\). Since the flow of
\(X=(y-\frac{1}{2\alpha})\frac{\partial}{\partial x}+\frac{\partial}{\partial y}\)
\[
  X_t(x,y)
  = \left(\begin{smallmatrix}1&t\\0&1\end{smallmatrix}\right)\left(\begin{smallmatrix}x\\y\end{smallmatrix}\right) + \left(\begin{smallmatrix}\frac{1}{2}t^2-\frac{1}{2\alpha}t\\t\end{smallmatrix}\right)
\] satisfies \(X_t^*\omega=\omega\), we see that \(X_t\) belongs to the
Lie algebra of \(\operatorname{Aut}({\mathcal{F}})\). Note that
\(X_{1/\alpha}=T\).

\end{rmenv}

The following result shows that this example is general.

\protect\phantomsection\label{corollary-1}
\begin{itenv}{Corollary 1}
If \(\operatorname{Aut}({\mathcal{F}})\) has an infinite number of
elements then the dimension of \(\operatorname{Aut}({\mathcal{F}})\) is
not zero.

\end{itenv}

\begin{proof}
Suppose, for contradiction, that
\(\dim\operatorname{Aut}({\mathcal{F}})=0\). Then, by
\hyperref[proposition-1]{Proposition 1},
\(\operatorname{Aut}({\mathcal{F}})\) has only a finite number of
elements, which contradicts the hypothesis.
\end{proof}

To obtain a holomorphic vector field in the Lie algebra of
\(\operatorname{Aut}({\mathcal{F}})\), as in
\hyperref[example-2]{Example 2}, we apply the Martinelli--Bochner
theorem \autocite{10} which says the following: the automorphism group
of a compact complex variety is a Lie group whose Lie algebra consists
of globally defined holomorphic vector fields on the variety.

\section{Automorphisms of holomorphic foliations on the projective
plane}\label{section-3}

\oldpage{52}

\protect\phantomsection\label{theorem-1}
\begin{itenv}{Theorem 1}
Let \({\mathcal{F}}\) be a holomorphic foliation on
\(\mathbb{P}_\mathbb{C}^2\). If
\(\#\operatorname{Aut}({\mathcal{F}})=+\infty\) in a suitable affine
coordinate system, then either \({\mathcal{F}}\) is induced by one of
the following \(1\)-forms:

\begin{enumerate}
\def\labelenumi{\arabic{enumi}.}
\tightlist
\item
  \(\omega=P\mathrm{d}{x}+(Q-yP)\mathrm{d}{y}\), where
  \(P,Q\in\mathbb{C}[x-\frac{y^2}{2}]\);
\item
  \(\omega=yp(y)\mathrm{d}{x}+(xp(y)-q(y))\mathrm{d}{y}\), where
  \(p,q\in\mathbb{C}[t]\);
\item
  \(\omega=yp(y^m/x^n)\mathrm{d}{x}+xq(y^m/x^n)\mathrm{d}{y}\)
\end{enumerate}

or \({\mathcal{F}}\) is linear.

\end{itenv}

\begin{proof}
Recall that the automorphism group of \(\mathbb{P}_\mathbb{C}^2\) is
isomorphic to \(\operatorname{PSL}(3,\mathbb{C})\). Since
\(\#\operatorname{Aut}({\mathcal{F}})=+\infty\),
\hyperref[corollary-1]{Corollary 1} tells us that
\(\operatorname{Aut}({\mathcal{F}})\) is a linear algebraic subgroup of
\(\operatorname{PSL}(3,\mathbb{C})\). In particular,
\(\operatorname{Aut}({\mathcal{F}})\) has a non-trivial Lie algebra, and
so, by the Martinelli--Bochner theorem, there exists a global vector
field \(X\) on \(\mathbb{P}_\mathbb{C}^2\) whose flow is in
\(\operatorname{Aut}({\mathcal{F}})\).

Any global holomorphic vector field \(X\) on \(\mathbb{P}_\mathbb{C}^2\)
can be represented in homogeneous coordinates \((x,y,z)\) in the form \[
  X(x,y,z)
  = M\cdot(x,y,z)^t
\] where \(M\in\operatorname{gl}(3,\mathbb{C})\). Analysing the possible
canonical Jordan forms, we have essentially three distinct cases.

Firstly, suppose that the canonical Jordan form of \(M\) is \[
  M
  = \begin{bmatrix}
    \alpha & 1 & 0
  \\0 & \alpha & 1
  \\0 & 0 & \alpha
  \end{bmatrix}
\] where \(\alpha\in\mathbb{C}^*\). In this case, the action \(\phi_t\)
induced by \(X\) can be written in the affine plane \(\{z=1\}\) as \[
  \phi_t(x,y)
  = (x+ty+\frac{t^2}{2}, y+t).
\] \oldpage{53} Therefore the orbits of \(X\) are rational curves, and
\(X\) admits a rational first integral \(x-y^2/2\). If \({\mathcal{F}}\)
is defined by the \(1\)-form
\(\omega=P(x,y)\mathrm{d}{x}+Q(x,y)\mathrm{d}{y}\) then
\(\phi_t^*\omega=\lambda(t)\omega\). Thus \(\lambda\equiv1\) and \[
  \begin{cases}
    P = P\circ\phi_t
  \\Q = tP\circ\phi_t + Q\circ\phi_t.
  \end{cases}
\] Thus \(\omega=P\mathrm{d}{x}+(R-yP)\mathrm{d}{y}\), where
\(P=P\circ\phi_t\) and \(R=R\circ\phi_t\). In other words, \(R\) and
\(P\) are integrals for the foliation defined by \(X\). Since
\(x-y^2/2\) is also a first integral and its fibres are connected, the
Stein factorisation theorem \autocite{8} implies that
\(P,R\in\mathbb{C}[x-y^2/2]\).

Now consider \(M\) of the form \[
  M
  = \begin{bmatrix}
    \alpha & 1 & 0
  \\0 & \alpha & 0
  \\0 & 0 & \beta
  \end{bmatrix}
\] where \(\alpha,\beta\in\mathbb{C}^*\). Here the induced action
\(\phi_t\) is, in affine coordinates \(\{z=1\}\), \[
  \phi_t(x,y)
  = e^{t(\alpha-\beta)}(x+ty,y).
\] As in the previous cases, \({\mathcal{F}}\) at \(\{z=1\}\) is defined
by \(\omega=P(x,y)\mathrm{d}{x}+Q(x,y)\mathrm{d}{y}\). Since
\(\phi_t^*\omega=\lambda(t)\omega\), we see that \[
  \lambda(t)P(x,y)
  = e^{t(\alpha-\beta)}P\circ\phi_t
  \tag{1}
\] \[
  \lambda(t)Q(x,y)
  = e^{t(\alpha-\beta)}(Q\circ\phi_t + tP\circ\phi_t).
  \tag{2}
\] Write \(P(x,y)\) as a sum of its homogeneous components
\(P(x,y)=\sum_j P_j(x,y)\). From (1) we deduce that, if \(P_j(x,y)\) is
not zero, then \(\lambda(t)=e^{t(\alpha-\beta)(j+1)}\). Comparing again
with (1), we see that \(P_j(x,y)=P_j(x+ty,y)\). This implies that
\(P_j\) does not depend on \(x\), and so \(P(x,y)=p_jy^j\). Using a
similar argument with \(Q(x,y)\), we can show that
\(Q(x,y)=-p_jxy^{j-1}+q_jy^j\). Then the \(1\)-form \(\omega\) can be
written as \[
  \omega
  = yp(y)\mathrm{d}{x} + (q(y)-xp(y))\mathrm{d}{y}.
\]

The remaining case is when \(M\) is diagonalisable, i.e. \[
  M
  = \begin{bmatrix}
    \alpha & 0 & 0
  \\0 & \beta & 0
  \\0 & 0 & \gamma
  \end{bmatrix}
\] where \(\alpha,\beta,\gamma\in\mathbb{C}^*\). \oldpage{54} As in the
previous cases, we can choose an affine coordinate system in which the
foliation \({\mathcal{F}}\) is described by
\(\omega=P\mathrm{d}{x}+Q\mathrm{d}{y}\) and the action \(\phi_t\)
induced by \(X\) is of the form \[
  \phi_t(x,y)
  = (e^{t(\alpha-\gamma)}x,e^{t(\beta-\gamma)y}).
\] Proceeding as before, we obtain the following relations: \[
  \lambda(t)P(x,y)
  = e^{t(\alpha-\gamma)}P\circ\phi_t
  \tag{3}
\] \[
  \lambda(t)Q(x,y)
  = e^{t(\beta-\gamma)}Q\circ\phi_t.
  \tag{4}
\] If we write \(P(x,y)=\sum_{i,j}p_{ij}x^iy^j\) and
\(Q(x,y)=\sum_{i,j}q_{ij}x^iy^j\) then we can deduce from (3) the
relations \[
  \lambda(t)
  = e^{t(\alpha-\gamma)(i+1)+t(\beta-\gamma)j}
  \quad\text{if }p_{ij}\neq0
  \tag{5}
\] \[
  \lambda(t)
  = e^{t(\alpha-\gamma)i+t(\beta-\gamma)(j+1)}
  \quad\text{if }q_{ij}\neq0.
  \tag{6}
\] Note that \(X\) admits a rational first integral if and only if
\((\alpha-\gamma)/(\beta-\gamma)=m/n\in\mathbb{Q}\).

Suppose that \(X\) has an rational first integral, and consider the
group \[
  G
  = \{ (l,k)\in\mathbb{Z}\oplus\mathbb{Z}\mid e^{t(\alpha-\gamma)l}=e^{t(\beta-\gamma)k}\text{ for all }t\in\mathbb{C}\}.
\] Since \(G\) is cyclic and generated by \((n,m)\), let
\((l_0,k_0)\in\mathbb{Z}\oplus\mathbb{Z}\) be such that
\(\lambda(t)=e^{t(\alpha-\gamma)l_0+t(\beta-\gamma)k_0}\). Then \[
  \omega
  = x^{l_0}y^{k_0}(yp(y^m/x^n)\mathrm{d}{x} + xq(y^m/x^n)\mathrm{d}{y}).
\]

Finally, assume that
\((\alpha-\gamma)/(\beta-\gamma)\not\in\mathbb{Q}\). Then there exists
\(t_0\in\mathbb{C}^*\) such that
\(e^{t_0(\alpha-\gamma)n}\neq e^{t_0(\beta-\gamma)m}\) for any integers
\(n\) and \(m\) that are not both zero. This implies that there exists a
unique pair of integers \((k,l)\) such that
\(\lambda(t_0)=e^{t_0(\alpha-\gamma)l+t_0(\beta-\gamma)k}\). As a
consequence of (3), we conclude that \(P(x,y)=p_{lk}x^{l-1}y^k\) and
\(Q(x,y)=q_{lk}x^ly^{k-1}\). Thus \[
  \omega
  = x^{l-1}y^{k-1}(p_{lk}y\mathrm{d}{x} + q_{lk}x\mathrm{d}{y})
\] and clearly \({\mathcal{F}}\) is a linear foliation.
\end{proof}

\oldpage{55}

\protect\phantomsection\label{corollary-2}
\begin{itenv}{Corollary 2}
Let \({\mathcal{F}}\) be a holomorphic foliation defined on
\(\mathbb{P}_\mathbb{C}^2\). Suppose that
\(\#\operatorname{Aut}({\mathcal{F}})=+\infty\). Then \({\mathcal{F}}\)
is birationally equivalent to a Riccati foliation.

\end{itenv}

\begin{proof}
From the proof of the \hyperref[theorem-1]{previous theorem} we have
that the foliation defined by \(X\) has no rational first integral, and
so \({\mathcal{F}}\) is linear. In particular, \({\mathcal{F}}\) is a
Riccati foliation after a blow-up.

Now suppose that the foliation \({\mathcal{G}}\) defined by \(X\) has a
first integral, i.e.~is a pencil of rational curves. Let
\(\sigma\colon M\to\mathbb{P}_\mathbb{C}^2\) be the minimal resolution
of \({\mathcal{G}}\). As we blow-up the singular points of \(X\), we
obtain a holomorphic field of vectors of \(\widetilde{X}\) defined on
all of \(M\) that belong to the Lie algebra of
\(\operatorname{Aut}(\sigma^*{\mathcal{F}})\).

Thus the geometric place of the tangencies of \(\sigma^*{\mathcal{F}}\)
and \(\widetilde{X}\) must be invariant under \(\sigma^*{\mathcal{F}}\),
see \autocite{6} and \autocite{7}. So \(\sigma^*{\mathcal{G}}\) is a
rational fibration, and we thus obtain that \(\sigma^*{\mathcal{F}}\) is
a Riccati foliation.
\end{proof}

\protect\phantomsection\label{proposition-2}
\begin{itenv}{Proposition 2}
Let \({\mathcal{F}}\) be a holomorphic foliation of general type on a
compact complex surface \(M\). Then
\(\operatorname{Aut}({\mathcal{F}})\) is isomorphic to a linear
algebraic group.

\end{itenv}

\begin{proof}
Since \({\mathcal{F}}\) is of general type, for a sufficiently large
integer \(m\), the \(m\)-th pluricanonical map \(\phi_m\) is a
bimeromorphism between \(M\) and the closure of its image, denoted by
\(N\).

Note that the group \(\operatorname{Aut}({\mathcal{F}})\) acts naturally
on the projectivisation of
\(\mathrm{H}^0(M,K_{\mathcal{F}}^{\otimes m})\). If \(\sigma\) is a
section of \(K_{\mathcal{F}}^{\otimes m}\) and \(\alpha\) is an
automorphism of \({\mathcal{F}}\), then the action is given by
\(\alpha(\sigma)=\alpha^*\sigma\).

Since \(\phi_m\) is a bimeromorphism between \(M\) and \(N\), this
action induces a monomorphism of groups \[
  \psi\colon \operatorname{Aut}({\mathcal{F}}) \to \operatorname{PSL}(k,\mathbb{C})
\] where
\(k=\dim_\mathbb{C}\mathrm{H}^0(M,K_{\mathcal{F}}^{\otimes m})-1\).

\oldpage{56}

Note that the image of \(\psi\) is precisely the automorphism of
\(\mathbb{P}_\mathbb{C}^{k-1}\) that leaves \(N\) and \({\mathcal{G}}\)
invariant, and so we can conclude that
\(\operatorname{Aut}({\mathcal{F}})\cong\operatorname{Aut}({\mathcal{G}})\)
is a linear algebraic subgroup of \(\operatorname{PSL}(k,\mathbb{C})\).
\end{proof}

In the proof of our principal result, we approximately follows the
arguments of Brunella in \autocite[86--88]{2}.

\begin{proof}
\emph{(Proof of \hyperref[theorem-A]{Theorem A}).} If
\(\#\operatorname{Aut}({\mathcal{F}})=+\infty\), then
\hyperref[proposition-2]{Proposition 2} implies that
\(\dim\operatorname{Aut}({\mathcal{F}}>0\), and by the
Martinelli--Bochner theorem we have a holomorphic vector field \(X\)
that determines a holomorphic foliation \({\mathcal{G}}\). Since on a
rational surface there exists a finite number of exceptional curves,
these have to be invariant under \(X\). Thus we can suppose that \(M\)
is minimal after making some contractions. If \(M\) is not
\(\mathbb{P}_\mathbb{C}^2\) then it is a Hirzebruch surface
\(\Sigma_n\). We first prove that there exists a singularity of the
foliation \({\mathcal{G}}\) outside the rational curve \(\Gamma_n\),
where \(\Gamma_n\) is the unique curve on \(\Sigma_n\) with the property
\(\Gamma^2=-n\). It is thus invariant under the flow of \(X\) and there
are two cases:

\begin{enumerate}
\def\labelenumi{\arabic{enumi}.}
\item
  \(X|_{\Gamma_n}=0\). If there is a fibre \(F\) (of a rational
  fibration of \(M\) on \(\Gamma_n\)) that is not invariant then we have
  that \(\operatorname{tang}({\mathcal{G}},F)>0\). Since
  \({\mathcal{G}}\) only has zero divisors,
  \(T_{\mathcal{G}}\cdot F\geqslant 0\), but on the other hand
  \(T_{\mathcal{G}}\cdot F=F^2-\operatorname{tang}({\mathcal{G}},F)=-\operatorname{tang}({\mathcal{G}},F)<0\).
  So \(F\) is invariant under \({\mathcal{G}}\), and so \(X\) is tangent
  to a rational fibration. Since
  \(\operatorname{tang}({\mathcal{G}},F)\) is invariant under
  \({\mathcal{F}}\) and \(X\) (see \autocite{6}), it must be a rational
  fibration or a Riccati foliation.
\item
  \(X|_{\Gamma_n}\neq0\). Then \(X\) has more than two zeros on
  \(\Gamma_n\) and, by the argument used in the previous case, the
  fibres that pass through these points are \({\mathcal{G}}\)-invariant.
  Furthermore, the Poincaré--Hopf index of \(X|_{\Gamma_n}\) agrees with
  the multiplicity of zeros of \(X|_{\Gamma_n}\). So we have two cases:

  \begin{enumerate}
  \def\labelenumii{\alph{enumii}.}
  \tightlist
  \item
    There are two \({\mathcal{G}}\)-invariant fibres, and there does not
    exist a singularity outside of \(\Gamma_n\), and so both
    singularities on \(\Gamma_n\) are saddle nodes and, by the index
    theorem, \(\Gamma_n^2=0\), which is a contradiction.
  \item
    If there is only one \({\mathcal{G}}\)-invariant fibre \(F\), then
    the multiplicity of \(X/_{\Gamma_n}\) at the singular point is
    \(2\). Now, if \(X\) has no singularities on \(F\) outside of
    \(\Gamma_n\), then \(X/_F\) has multiplicity \(2\). Thus, again by
    the index theorem, \(\Gamma_n^2=0\), which is a contradiction.
  \end{enumerate}
\end{enumerate}

Resuming, we have a fibre \(F\) and a point \(p\in F\) such that
\(X(p)=0\) and \(p\not\in\Gamma_n\). Blowing-up at \(p\) and
blowing-down the strict transform of \(F\) we arrive at
\(\Sigma_{n-1}\). Since \(X\) is still holomorphic, we blow-up at
singularities of \(X\) and shrink invariant curves. Following this
procedure we reduce our problem to
\(\Sigma_1=\mathbb{P}_\mathbb{C}^1\times\mathbb{P}_\mathbb{C}^1\), and
in this case \(X=X_1\oplus X_2\), where \(X_i\) is a holomorphic vector
field on \(\mathbb{P}_\mathbb{C}^1\).

From here we apply a blow-up and two blow-downs to arrive at
\(\mathbb{P}_\mathbb{C}^2\), with \(X\) still a holomorphic vector
field.\footnote{\emph{{[}Trans.{]} The original has a figure here,
  showing this bimeromorphism between
  \(\mathbb{P}_\mathbb{C}^1\times\mathbb{P}_\mathbb{C}^1\) and
  \(\mathbb{P}_\mathbb{C}^2\), which I have not tried to reproduce.}} To
finish, we apply \hyperref[corollary-2]{Corollary 2}.
\end{proof}

\begin{proof}
\emph{(Proof of \hyperref[theorem-B]{Theorem B}).} By Theorem 3.3.1 in
\autocite{12}, we know that the Riccati foliations and rational
fibrations have Kodaira dimension less than \(2\). Thus by
\hyperref[theorem-A]{Theorem A} we know that, if a foliation is of
general type, then its automorphism group must be finite.
\end{proof}

\subsection{Final commentary}\label{section-4}

This result that we have proven in this article for rational surfaces is
valid for any complex surface, see \autocite{6} and \autocite{7}. In
this latter article we also prove some results in dimension higher than
\(2\), on the characterisation of foliations with infinite automorphism
groups that are algebraic linear. Some questions naturally arise, such
as ``what is the optimal dimension for the group
\(\operatorname{Aut}({\mathcal{F}})\) of a foliation of general type?'',
but were not answered due to the difficulties that appear. This question
is relevant since the unique foliation with known automorphism group is
the Jouanoulo foliation, see \autocite{9}.

\subsection{Thanks}\label{section-5}

This article is a summary of a part of my doctoral thesis \autocite{6}
which was written under the direction of Professor Cesar Camacho and an
article co-authored with Jorge Pereira, both of whom I thank for the
many comments and suggestions.


% Bibliography
\nocite{*}
\printbibliography[heading=bibintoc,title=Bibliography]

\end{document}
